\documentclass[12pt, twoside, openright]{report} % Fuente a 12pt, formato doble página y chapter a la derecha
\raggedbottom % No ajustar el contenido con un salto de página

% MÁRGENES: 2,5 cm sup. e inf.; 3 cm izdo. y dcho.
\usepackage[
a4paper,
vmargin=2.5cm,
hmargin=3cm
]{geometry}

% INTERLINEADO: Estrecho (6 ptos./interlineado 1,15) o Moderado (6 ptos./interlineado 1,5)
\renewcommand{\baselinestretch}{1.15}
\parskip=6pt

% DEFINICIÓN DE COLORES para portada y listados de código
\usepackage[table]{xcolor}
\definecolor{azulUC3M}{RGB}{0,0,102}
\definecolor{gray97}{gray}{.97}
\definecolor{gray75}{gray}{.75}
\definecolor{gray45}{gray}{.45}

% Soporte para GENERAR PDF/A
\usepackage{etoolbox}
\makeatletter
\@ifl@t@r\fmtversion{2021-06-01}%
 {\AddToHook{package/after/xmpincl}
   {\patchcmd\mcs@xmpincl@patchFile{\if\par}{\ifx\par}{}{\fail}}}{}
\makeatother
\usepackage[a-1b]{pdfx}

% ENLACES
\usepackage{hyperref}
\hypersetup{colorlinks=true,
  linkcolor=black, % enlaces a partes del documento (p.e. índice) en color negro
  urlcolor=blue} % enlaces a recursos fuera del documento en azul

% Añadir pdfs como partes del documento
\usepackage{pdfpages}

% Quitar la indentación de principio de los párrafos
\setlength{\parindent}{0em}

% EXPRESIONES MATEMÁTICAS
\usepackage{amsmath,amssymb,amsfonts,amsthm}

\usepackage{txfonts} 
\usepackage[T1]{fontenc}
\usepackage[utf8]{inputenc}

% Insertar gráficas y fotos
\usepackage{tikz}
\usetikzlibrary{cd}
\usepackage{pgfplots}

\usepackage[spanish, es-tabla]{babel} 
\usepackage[babel, spanish=spanish]{csquotes}
\AtBeginEnvironment{quote}{\small}

% diseño de PIE DE PÁGINA
\usepackage{fancyhdr}
\pagestyle{fancy}
\fancyhf{}
\renewcommand{\headrulewidth}{0pt}
\fancyfoot[LE,RO]{\thepage}
\fancypagestyle{plain}{\pagestyle{fancy}}

% DISEÑO DE LOS TÍTULOS de las partes del trabajo (capítulos y epígrafes o subcapítulos)
\usepackage{titlesec}
\usepackage{titletoc}
\titleformat{\chapter}[block]
{\large\bfseries\filcenter}
{\thechapter.}
{5pt}
{\MakeUppercase}
{}
\titlespacing{\chapter}{0pt}{0pt}{*3}
\titlecontents{chapter}
[0pt]                                               
{}
{\contentsmargin{0pt}\thecontentslabel.\enspace\uppercase}
{\contentsmargin{0pt}\uppercase}                        
{\titlerule*[.7pc]{.}\contentspage}                 

\titleformat{\section}
{\bfseries}
{\thesection.}
{5pt}
{}
\titlecontents{section}
[5pt]                                               
{}
{\contentsmargin{0pt}\thecontentslabel.\enspace}
{\contentsmargin{0pt}}
{\titlerule*[.7pc]{.}\contentspage}

\titleformat{\subsection}
{\normalsize\bfseries}
{\thesubsection.}
{5pt}
{}
\titlecontents{subsection}
[10pt]                                               
{}
{\contentsmargin{0pt}                          
  \thecontentslabel.\enspace}
{\contentsmargin{0pt}}                        
{\titlerule*[.7pc]{.}\contentspage}  

\usepackage{adjustbox}

% DISEÑO DE TABLAS.
\usepackage{multirow} % permite combinar celdas 
\usepackage{caption} % para personalizar el título de tablas y figuras
\usepackage{floatrow} % utilizamos este paquete y sus macros \ttabbox y \ffigbox para alinear los nombres de tablas y figuras de acuerdo con el estilo definido. Para su uso ver archivo de ejemplo 
\usepackage{array} % con este paquete podemos definir en la siguiente línea un nuevo tipo de columna para tablas: ancho personalizado y contenido centrado
\newcolumntype{P}[1]{>{\centering\arraybackslash}p{#1}}
\DeclareCaptionFormat{upper}{#1#2\uppercase{#3}\par}

% Diseño de tabla para ingeniería
\captionsetup[table]{
  format=hang,
  name=Tabla,
  justification=centering,
  labelsep=colon,
  width=.75\linewidth,
  labelfont=small,
  font=small,
}

% DISEÑO DE FIGURAS.
\usepackage{graphicx}
\graphicspath{{img/}} %ruta a la carpeta de imágenes

% Diseño de figuras para ingeniería
\captionsetup[figure]{
  format=hang,
  name=Fig.,
  singlelinecheck=off,
  labelsep=colon,
  labelfont=small,
  font=small    
}

% NOTAS A PIE DE PÁGINA
\usepackage{chngcntr} % Para numeración continua de las notas al pie
\counterwithout{footnote}{chapter}

% LISTADOS DE CÓDIGO
% soporte y estilo para listados de código. Más información en https://es.wikibooks.org/wiki/Manual_de_LaTeX/Listados_de_código/Listados_con_listings
\usepackage{listings}

% definimos un estilo de listings
\lstdefinestyle{estilo}{ frame=Ltb,
  framerule=0pt,
  aboveskip=0.5cm,
  framextopmargin=3pt,
  framexbottommargin=3pt,
  framexleftmargin=0.4cm,
  framesep=0pt,
  rulesep=.4pt,
  backgroundcolor=\color{gray97},
  rulesepcolor=\color{black},
  %
  basicstyle=\ttfamily\footnotesize,
  keywordstyle=\bfseries,
  stringstyle=\ttfamily,
  showstringspaces = false,
  commentstyle=\color{gray45},     
  %
  numbers=left,
  numbersep=15pt,
  numberstyle=\tiny,
  numberfirstline = false,
  breaklines=true,
  xleftmargin=\parindent
}

\captionsetup[lstlisting]{font=small, labelsep=period}
% fijamos el estilo a utilizar 
\lstset{style=estilo}
\renewcommand{\lstlistingname}{\uppercase{Código}}

\pgfplotsset{compat=1.17} 
%-------------
% DOCUMENTO
%-------------

\begin{document}
\pagenumbering{roman} % Se utilizan cifras romanas en la numeración de las páginas previas al cuerpo del trabajo

%----------
% PORTADA
%---------- 
\begin{titlepage}
	\begin{sffamily}
		\color{azulUC3M}
		\begin{center}
			\begin{figure}[H] % Incluimos el logotipo de la Universidad
				\makebox[\textwidth][c]{\includegraphics[width=16cm]{Portada_Logo.png}}
			\end{figure}
			\vspace{2.5cm}
			\begin{Large}
				Grado en Ingeniería Informática\\
				2021-2022\\
				\vspace{2cm}
				\textsl{Apuntes}\\
				\bigskip
			\end{Large}
			{\Huge Teoría avanzada de la computación}\\
			\vspace*{0.5cm}
			\rule{10.5cm}{0.1mm}\\
			\vspace*{0.9cm}
			{\LARGE Jorge Rodríguez Fraile\footnote{\href{mailto:100405951@alumnos.uc3m.es}{Universidad: 100405951@alumnos.uc3m.es}  |  \href{mailto:jrf1616@gmail.com}{Personal: jrf1616@gmail.com}}}\\
			\vspace*{1cm}
		\end{center}
		\vfill
		\color{black}
		\includegraphics[width=4.2cm]{img/creativecommons.png}\\
		Esta obra se encuentra sujeta a la licencia Creative Commons\\ \textbf{Reconocimiento - No Comercial - Sin Obra Derivada}
	\end{sffamily}
\end{titlepage}

%----------
% ÍNDICES
%---------- 

%--
% Índice general
%-
\tableofcontents
\thispagestyle{fancy}

%--
% Índice de figuras. Si no se incluyen, comenta las líneas siguientes
%-
\listoffigures
\thispagestyle{fancy}

%--
% Índice de tablas. Si no se incluyen, comenta las líneas siguientes
%-
\listoftables
\thispagestyle{fancy}

%----------
% TRABAJO
%---------- 

\pagenumbering{arabic} % numeración con números arábigos para el resto de la publicación  

%----------
% COMENZAR A ESCRIBIR AQUÍ
%---------- 

\chapter{Información}\label{ch:informacion}
\section{Profesores}\label{sec:profesores}
\begin{quote}
	Magistral: Jose Luis Mira Peidro, jmira@inf.uc3m.es

	Prácticas: Juan Manuel Alonso Weber, jmaw@ia.uc3m.es
\end{quote}

\section{Sistema de evaluación}\label{sec:sistema-de-evaluación}
\begin{itemize}
	\item 34 \% Examen Final (nota mínima 4/10)
	\item 66 \% Evaluación continua (EC)
	      \begin{itemize}
		      \item 11 \% Práctica 1
		      \item 11 \% Práctica 2
		      \item 11 \% Práctica 3
		      \item 33 \% Examen sobre las prácticas
	      \end{itemize}
\end{itemize}

\section{Examen Practicas}\label{sec:examen-practicas}
Para el examen final no entra la recursividad no homogénea.

Se pregunta sobre las 2 primeras prácticas

De AKS se pregunta del material que se entregó, para verificar que se sabe de todo.

Mirar además de la práctica, el material y las cuestiones.
\begin{itemize}
	\item Ejem. Describir brevemente el algoritmo.
	\item Ejem. Saber de la condición suficiente.
\end{itemize}

Se puede preguntar en AKS sobre alguna variación sobre el problema, para saber cómo se comportaría, de manera que se sabe cómo se procedió en el original.

Habrá preguntas más de tipo test y preguntas de desarrollar.

Ventajas de AKS sobre otros algoritmos.

No entra la tercera práctica.

\chapter{Análisis de la complejidad de los Algoritmos}\label{ch:análisis-de-la-complejidad-de-los-algoritmos}
\adjustbox{width=\textwidth,center}{%
  \begin{tikzcd}
                                      &                                                             & \text{Complejidad Computacional} \arrow[d] &                                       \\
                                      &                                                             & \text{tipos} \arrow[rd] \arrow[ld]         &                                       \\
                                      & \text{Complejidad temporal} \arrow[d]                       &                                            & \text{Complejidad espacial} \arrow[d] \\
                                      & S(n) \arrow[d]                                              &                                            & T(n)                                  \\
                                      & \text{Número de pasos base} \arrow[ld] \arrow[d] \arrow[rd] &                                            &                                       \\
    \text{Máquina de Turing} \arrow[d] & \text{Ensamblador} \arrow[d]                                & \text{Lenguaje alto nivel*} \arrow[d]      &                                       \\
    \text{transición}                  & \text{instrucción}                                          & \text{operación elemental} \arrow[d]       &                                       \\
                                      &                                                             & \text{OE}                                  &
  \end{tikzcd}
}

* Estructuras de los algoritmos siempre se pueden considerar como estructuras de tipo Secuencial, Iterativo y Condicional, esto determina la operación elemental.

\section{Introducción}\label{sec:introduccion}
\textbf{Complejidad Computacional:} Recursos requeridos durante la ejecución de un algoritmo que da respuesta a un problema. Hay 2 tipos:
\begin{itemize}
  \item Tiempo: Número de pasos base de ejecución de un algoritmo para resolver un problema.
  \item Espacio: Memoria utilizada para resolver un problema. Registros, Memoria caché y Memoria RAM, etc.
\end{itemize}
Espacial y Temporal.

\textbf{Problema:} Una función entre el espacio de instancias y el espacio de respuestas, $X \rightarrow P(X)$.
\begin{itemize}
  \item Ejem. La ecuación de grado 2 tiene raíces cuadradas? $P(x^2+1=0)=\text{No}$
\end{itemize}

\textbf{Problema de decisión:} Tiene 2 respuestas: Si o No.
\pagebreak

\textbf{Instancia de un problema:} Es la especificación exacta de los datos de un problema para un caso particular.
\begin{itemize}
  \item \textbf{Tamaño de instancia:} Número de datos de la instancia.
  \item Ejem. Independiente, $x^2$ y $x$ son los 3 parámetros de una instancia de un problema de ecuación de segundo grado.
\end{itemize}

\textbf{Algoritmo:} Conjunto de instrucciones que garantiza encontrar una solución correcta para cualquier instancia en un número finito de pasos.

La complejidad puede hacer referencia a un algoritmo o un problema.
\begin{itemize}
  \item \textbf{Algoritmo:} Resolver las instancias de un problema. \\ Medida de número de pasos base requeridos para la PEOR instancia de tamaño n. \\ Relación entre n y el número de pasos.
  \item \textbf{Problema:} Es el objeto de estudio de la Teoría de la Complejidad Computacional.
\end{itemize}

\section{Paso Base}\label{sec:paso-base}
Se pueden modelar los algoritmos mediante:
\begin{enumerate}
    \item Construcción matemática denominada Máquina de Turing $\rightarrow$ Transición / Tabla de transiciones.
  \item Código Ensamblador $\rightarrow$ Instrucción \ Conjunto estructurado de instrucciones.
  \item Lenguaje de alto nivel $\rightarrow$ Operación Elemental / Conjunto estructurado de instrucciones.
    \begin{itemize}
      \item Suma, Restar División, Multiplicación, etc.
    \end{itemize}
\end{enumerate}

\section{Complejidad de los Algoritmos}\label{sec:complejidad-de-los-algoritmos}
Se expresa como: $T(n)$, número de pasos peor caso con entrada n.

La complejidad computacional puede referirse también a los recursos de espacio necesarios según el tamaño de la entrada. Así tenemos que:
\begin{itemize}
  \item Tiempo de ejecución (complejidad temporal) $T(n)$
  \item Espacio necesario para los datos (complejidad espacial) $S(n)$
\end{itemize}

Todo algoritmo puede ser implementado con 3 estructuras: Secuencial, Condicional e Iterativa.

\textbf{Operación Elemental (OE)}
\begin{itemize}
  \item Comparación
  \item Asignación variable
  \item Acceso a estructura básica
  \item Llamada función
  \item Retorno función
\end{itemize}

\textbf{Estructura Secuencial:} Cada I es una instancia o bloque, se calcula:
\begin{itemize}
  \item El de todos: $T(I_1; I_2; \dots; I_n) = T(I_1) + T(I_2) + \dots + T(I_n)$
  \item General: $T= \max \{T(I_1), T(I_2), \dots, T(I_n)\}$ El T es el más alto, el del peor caso/bloque.
  \item Para los programas es la suma de todas las estructuras: $T = T_1 + T_2 + \dots + T_k = \max \{ T_1, T_2, \dots, T_k \}$
\end{itemize}

\textbf{Estructura Condicional:}
\begin{itemize}
  \item if (Condición) then BloqueThen else BloqueElse
  \item $T = T(Condicion) + \max \{T(BloqueThen), T(BloqueElse)\}$
\end{itemize}

\textbf{Estructura Iterativa:}
\begin{itemize}
  \item while(C) S - $T = C + n_{iter} (C+S)$
  \item for(A C B) S - $T = A + C + n_{iter} (C+B+S)$
  \begin{itemize}
    \item A: Inicialización
    \item C: Condición
    \item B: Incremento
    \item S: Bloque
  \end{itemize}
\end{itemize}
\pagebreak

\textbf{Llamada a Procedimiento, Función o Método:}
\begin{itemize}
  \item $funcion(p_1, p_2, \dots, p_k) f$
  \item 1 (por la llamada) + tiempo de evaluación parámetros + tiempo que tarde en ejecutarse el ''cuerpo'' de la función.
  \item $T = 1 + T(p_1) + T(p_2) + \dots + T(p_n) + T(f)$
\end{itemize}

No se contabiliza la copia de los argumentos a la pila de ejecución salvo estructuras complejas (vector, registros) que se pasen por valor.

El paso por referencia y paso de punteros, no se contabiliza.

\subsection{Ejemplo}\label{subsec:ejemplo}
\begin{lstlisting}[language=Java,label={lst:lstlisting}]
for (int i = 0, acum = 0; i < n; i++)
  acum += i;
\end{lstlisting}

$T(n) = A + C + n ( C + B + S ) = 2 + 1 + n ( 1 + 2 + 2 ) = 3 + 5 n$

\begin{lstlisting}[language=Java,label={lst:lstlisting2}]
if (x < 5)
  for (int i = 0; i < n; i++)
    acum += i;
else
  acum = 3;
\end{lstlisting}

$T(n) = 1 + \max \{ 1 + 1 + n(1+2+2) ; 1 \} = 1 + \max \{ 2 + 5 n ; 1 \} = 3 + 5n$

\begin{lstlisting}[language=Java,label={lst:lstlisting3}]
for (int i = 0, acum = 0; i < n; i++)
  for (int j = 0; j < i/2; j++)
    acum += i;
\end{lstlisting}

$T(n) = 2 + 1 + n ( 1 + 2 + [ 1 + 2 + \frac{n}{2} ( 2 + 2 + 2) ]) = 3 + n (3 + 3 + 3 n) = 3 + 6 n + 3 n^2$

\begin{lstlisting}[language=Java,label={lst:lstlisting4}]
x = 5;
while (x < N)
  x = 3 * x;
\end{lstlisting}

$T(n) = 1 + 1 + p ( 1 + 2 ) = 2 + 3 p$

$5 \cdot 3^p < n; p \log 3 < \log \frac{n}{7}; p = \frac{\log \frac{n}{5}}{\log 3}$
\pagebreak

\section{Algoritmos Recursivos}\label{sec:algoritmos-recursivos}
\adjustbox{width=\textwidth,center}{%
\begin{tikzcd}
  & \text{Algoritmo recursivo } \text{(Tiene coste } T(n) \text{)} \arrow[d] &                           \\
  & \text{Ecuaciones de recurrencia } E(n) \arrow[ld] \arrow[rd]     &                           \\
\text{Ecuacion caracteristica (Polinomica } P(n)=0 \text{)} \arrow[d] &                                                                  & \text{Cambio de variable} \\
{\text{Obtenemos } x_1, x_2, x_3, ..., x_n} \arrow[d]                 &                                                                  &                           \\
\text{Obtenemos de } T(n) \text{ no recursiva}                        &                                                                  &                          
\end{tikzcd}
}

El coste ($T(n)$) de un algoritmo recursivo será también recursivo.

$T(n)= E(n)$ En la expresión 
\begin{itemize}
  \item $E$ aparece la propia función $T$.
  \item El objetivo es encontrar $T(n)= f(n)$, sin funciones $T$.
\end{itemize}

Resolver la ecuación de recurrencia. Dos métodos:
\begin{enumerate}
  \item Despliegue de la recurrencia
  \item Resolución general de recurrencias
  \begin{enumerate}
    \item Recurrencia homogénea (ecuación característica)
    \item Recurrencia útiles (fórmulas maestras)
    \item Recurrencias no homogéneas y no lineales (no se tratan en este curso)
  \end{enumerate}
  
\end{enumerate}
  
\subsection{Despliegue de Recurrencias}
Aplicar varias veces la fórmula recurrente hasta obtener una fórmula general que relaciona la función para el tamaño original con otros tamaños menores. A partir de esta fórmula se obtiene otra que la relaciona con el caso base.

Pasos:
\begin{enumerate}
  \item Ecuación ($T(n)=...$)
  \item Despliegue, ($T(n-1), T(n-2), etc$)
  \item Aplicaciones del caso base para $n=0, ..., k\;\; T(0)=a, T(1)=b, ...$
  \item Resolución
\end{enumerate}

No debe aplicarse este método si aparecen varios términos recurrentes, como es el caso de la serie de Fibonacci.

Se parte de un caso base y se va sustituyendo dentro de la ecuación el propio término, de manera que cuando vamos sustituyendo y resolviendo la ecuación hasta llegar a un caso en que podamos aplicar el caso base. En ese punto podremos sustituir nuestro caso base en la ecuación y quitarnos completamente la recurrencia, t().

Lo que buscamos es que al desdoblar los términos constantemente encontremos una regularidad. En ese momento hacemos que una variable valga un valor que podamos utilizar para usar el caso base, normalmente t(1).

Ejem. Factorial
\begin{itemize}
  \item $t(1) = 1 \textit{ y } t(n) = t(n-1)+1$
  \item $t(n) = t(n-2)+1+1=t(n-2)+2=t(n-3)+1+2=t(n-3)+3=t(n-k)+k$
  \item Sea $k=n-1$ $t(n)=t(n-(n-1))+n-1=t(1)+n-1=n$
\end{itemize}

\subsection{Resolución General de la Ecuación de Recurrencia}
$c_nt(n)+c_{n-1}t(n-1)+...+c_{n-k}t(n-k)=b$, expresión en n.

Si $b=0$, se dice que es Ecuación de Recurrencia Lineal (de orden k) HOMOGÉNEA

Los términos de las ecuaciones de recurrencia están afectados por coeficientes. Debemos encontrar una expresión no recursiva de $t$.

No nos ocuparemos de ecuaciones en las que el coeficiente de $t(n-k)$, i.e. $c_{n-k}$ es variable (depende de n), trataremos solo con: ECUACIONES DE RECURRENCIA LINEALES CON COEFICIENTES CONSTANTES.

\subsubsection{Ecuación de Recurrencia Lineal Homogenea}
Método de la Ecuación Característica.
\begin{itemize}
  \item A una ecuación de recurrencia de orden k, mediante un cambio de variable, se le asocia una ecuación polinómica de grado k (la ecuación característica).
  \item Las soluciones de la ecuación polinómica proporcionan las soluciones de la ecuación de recurrencia.
  \item Solución general de la ecuación de recurrencia:
  \begin{itemize}
    \item Es combinación lineal de las soluciones básicas.
    \item Las constantes de la combinación lineal se calculan imponiendo los casos base.
    \item Con ello obtenemos la solución (no recursiva) de la recurrencia $T(n)$
  \end{itemize}

\end{itemize}

Cambio de variable: $T(n)= x^n$, se obtiene la ecuación característica $c_n x^n+ c_{n-1}x^{n-1}+ ... + c_{n-k}x^{n-k} = 0$ siendo $r_1, r_2, ..., r_k$ las raíces reales o complejas.

A una ecuación de recurrencia de orden k, mediante un cambio de variable, se le asocia una ecuación polinómica de grado k (la ecuación característica).

Lo que hacemos tras la sustitución es obtener las raíces, que pueden ser distintas $x=2$ y $x=3$ o iguales $x=1$, $(x-1) \cdot (x-1)$.
\begin{itemize}
  \item Para las raíces distintas se aplica: \\ 
  $T(n) = \sum _1 ^k b_i r_i^n, \; b_i, r_i \in \mathbb{R}$ \\
  $r_1$ y $r_2$ son raíces diferentes, $T(n)=b_1r_1^n+b_2r_2^n$ es solución.
  \item Una raíz con multiplicidad mayor que 1: \\ 
  $T(n) = \sum _{i=1} ^m b_i n^{i-1}r_1^n+\sum _{i=m+1} ^k b_i r_{i-m+1}^n, \; b_i, r_i \in \mathbb{R}$ \\ 
  $r_1$ son las raíces de la ecuación característica y $m$ la multiplicidad cuando es mayor que 1.\\
  $r$  raíz doble, $T(n)=b_1nr^n+b_2r^n$ es solución.
\end{itemize}

Se cambia el $t(n)$, $t(n-1)$, \ldots por variables $x$ que dependan del $n$.

El grado de la variable siempre es un grado menos que el número de términos recurrentes y cada uno de valor menor reduce el grado.

Ejem. $T(n)=T(n-1)+T(n-1) \rightarrow x^2= x + 1$
\begin{itemize}
  \item Sus raíces son $r = \frac{1\pm \sqrt{5}}{2}$, estos términos se sustituirían en las ecuaciones de arriba quedando:
  \item Sabiendo $T(0) = 0$ y $T(1)=1$
  \begin{itemize}
    \item $T(n)= b_1(\frac{1+ \sqrt{5}}{2})^n + b_2(\frac{1- \sqrt{5}}{2})^n$
    \item $T(0)= b_1(\frac{1+ \sqrt{5}}{2})^0 + b_2(\frac{1- \sqrt{5}}{2})^0 = b_1 + b_2$
    \item $T(1)= b_1(\frac{1+ \sqrt{5}}{2})^1 + b_2(\frac{1- \sqrt{5}}{2})^1 = 1$
  \end{itemize}
  \item Despejamos con ambas ecuaciones: $b_1 = -b_2 = \frac{1}{5}$.
\end{itemize}

\section{Trazabilidad}\label{sec:tratabilidad}


\section{Análisis Asintótico}\label{sec:análisis-asintótico}

\chapter{Teoría de la Computabilidad}\label{ch:teoria-de-la-computabilidad}
Clasificar los distintos problemas y formalizar el concepto de computar.

$\left\{\begin{matrix}
    Complejidad& \left\{\begin{matrix}
    Algoritmo &\rightarrow& \textit{Coste temporal } T(n)\\ 
    Problemas &\rightarrow& \textit{Tipo de problema según complejidad} \left\{\begin{matrix}
    P\\ 
    NP\\ 
    NPC
    \end{matrix}\right.
    \end{matrix}\right.\\ 
    Computabilidad& \rightarrow Problema \left\{\begin{matrix}
    Computable\\ 
    \textit{No computable}
    \end{matrix}\right.
    \end{matrix}\right.$

\textbf{Problema:} Descripción de una pregunta que requiere una respuesta.
\begin{itemize}
    \item p: instancia $\rightarrow$ respuesta
    \item $x \rightarrow p(x)$
\end{itemize}

\textbf{Instancia:} Especificación exacta de los datos de un problema para un caso particular.

\textbf{Algoritmo:} Conjunto de instrucciones que garantiza encontrar una solución correcta para cualquier instancia en un \textbf{número finito de pasos}.

\textbf{Analogía geométrica:} Hallar el lado de un cuadrado inscrito en un círculo, se puede dibujar fácilmente, pero en cuanto a matemáticas nunca se halla exactamente al tener raíces cuadradas. $r^2 = \frac{l^2}{2}+\frac{l^2}{2}; \; l= \sqrt{2}r$. Un caso en el que no es posible de ambas manera es dibujar un cuadrado con la misma área que un círculo, dado que entra en juego $\sqrt{\pi}$. $l^2 = \pi r^2; \; l = \sqrt{\pi}r$

\section{Problema de Satisfabilidad Booleana: SAT}
Un problema SAT2 puede ser...

if x then y else z

$(x \wedge y) \vee (\overline{x} \wedge z)$

\section{Clique}
Subgrafo completo, en el que todos los vértices están conectados con todos.

Tiene mucha potencia social, todos conocen o comparten gustos.

\section{Decidibilidad}
Un problema de computación puede considerarse como un lenguaje.

Las posibles soluciones del problema pueden considerarse como las palabras que pertenecen al lenguaje.
\begin{itemize}
    \item Ejem. $ax^2+bx+c=0$
    \item Problema: Tiene solución la ecuación en $\mathbb{R}$
    \item Respuestas $\in \{Si, No\}$
    \item Soluciones $\in \mathbb{R}$
    \item Lenguaje: $L= \{a,b,c\}$
\end{itemize}

\begin{figure}[H]
    \ffigbox[\FBwidth]
    {\caption{Diagrama Decidibilidad}}
    {\tikzset{every picture/.style={line width=0.75pt}}
    \begin{tikzpicture}[x=0.75pt,y=0.75pt,yscale=-1,xscale=1]
    \draw   (179,100) -- (380,100) -- (380,220) -- (179,220) -- cycle ;
    \draw   (210,160) .. controls (210,137.91) and (227.91,120) .. (250,120) .. controls (272.09,120) and (290,137.91) .. (290,160) .. controls (290,182.09) and (272.09,200) .. (250,200) .. controls (227.91,200) and (210,182.09) .. (210,160) -- cycle ;
    \draw    (320,100) -- (320,220) ;
    
    \draw (181,103) node [anchor=north west][inner sep=0.75pt]   [align=left] {T-R};
    \draw (321,102) node [anchor=north west][inner sep=0.75pt]   [align=left] {T-N-R};
    \draw (237.5,152) node [anchor=north west][inner sep=0.75pt]   [align=left] {T-D};
    \end{tikzpicture}}
\end{figure}
\begin{itemize}
    \item \textbf{T-D:} Turing-Decidible. Es un problema para el cual una MT acepta todas las soluciones y rechaza las que no son. \\ Se para siempre.
    \item \textbf{T-R:} Turing-Reconocible, pero no decidible. Es un problema para el cual la MT acepta todas las que son soluciones. \\ Puede entrar en bucle para las no soluciones.
    \item \textbf{T-N-R:} Turing-No-Reconocible. Ni acepta las que son soluciones, ni rechaza las que no son. \\ Entra en bucle.
\end{itemize}

Un T-D es un T-R, siguiendo el diagrama

\textbf{Reconocer} un problema (lenguaje) es aceptar todas las soluciones del problema.

\textbf{Decidir} un problema es:
\begin{itemize}
    \item Reconocer todas las soluciones.
    \item Rechazar todas las no-soluciones.
\end{itemize}

Solo los T-D se llaman Computables o Decidibles.
\begin{itemize}
    \item Ejem. El teorema de Fermat ($x^n + y^n = z^n$ con $n>2 / n \in \mathbb{R}$) es indecible porque ni computacional ni matemáticamente no se puede demostrar, aunque se sabe que no hay otro real que cumpla la condición aparte de 2. La máquina estaría infinitamente probando valores.
\end{itemize}

CUIDADO. Indecidible no especifica si es T-R o T-N-R, a veces al T-R se le llama parcialmente decidible.

\textbf{MT Reconocedora:} Que no genera salida lee la cinta, pero no lo modifica/manipula.

\textbf{MT Transductora:} Que genera salida. Modifica la cinta, además de reconocer la cinta.

\chapter{Teoría de la Complejidad Computacional}\label{ch:teoría-de-la-complejidad-computacional}
Consideremos 2 tipos de recursos:
\begin{itemize}
  \item Tiempo: tiempo de ejecución.
  \item Espacio: recursos de memoria utilización.
\end{itemize}

\section{La Máquina de Turing y medida de la Complejidad}
El tiempo de ejecución de Máquina de Turing es la función $T(n): N \rightarrow N$, donde T(n) es el número máximo de pasos que una MT utiliza para realizar el cómputo para entrada de tamaño n.

Complejidad Computacional: se refiere a la complejidad de los problemas
\begin{itemize}
  \item Se toma la peor instancia
  \item Se utiliza la cota Superior  Asintótica: O(g)
\end{itemize}

La principal herramienta para comparar problemas y clasificarlos:
\begin{itemize}
  \item Desde el punto de vista de la computabilidad
  \item Desde el punto de vista de la complejidad computacional
\end{itemize} 

\section{Reducción entre problemas}
Dado dos Lenguajes A y B, y una función $f: \Sigma^* \rightarrow \Sigma^*$.

f es una transformación del lenguaje universal del problema A al lenguaje universal del problema B. Se dice que \textbf{A es reducible a B}, bajo f si:

$$\begin{matrix}
  \forall w \in A \Leftrightarrow f(w)\in B \\ 
  A \leq_f B
\end{matrix}$$

\textbf{Reducción:} transformar un problema en otro. Transformar toda instancia de un problema en instancias de otro problema.

Intuitivamente
\begin{itemize}
  \item es una manera de convertir un problema a otro problema, de tal manera que la solución al segundo problema se puede utilizar para resolver el primer problema.
  \item ''Si el problema A es reducible al problema B, entonces la solución al problema B se
  puede utilizar para resolver el problema A''
\end{itemize}

Formalmente
\begin{itemize}
  \item Si el coste de Reducir A en B, es decir, el proceso de transformar una instancia de A en una instancia de B, es de orden polinómico (o menor) se puede afirmar que:
  \begin{itemize}
    \item Resolver A no puede ser más complejo (''harder'') que resolver B. Es decir, resolver A a lo sumo es tan complejo de resolver como B.
    \item Resolver B al menos es tan complejo como resolver A
    \item Estas dos afirmaciones se expresan así: $A \leq B$
  \end{itemize}
  \item Reducir el problema A al problema B en tiempo polinómico implica que la reducción se lleva a cabo con una función cuyo $T(n)=O(n^k)$, $A \leq_p B$
  \item Si $A \leq_p B$ y $B \leq_p A$ son igual de difíciles.
\end{itemize}

\section{Clases de Problemas}
Tiempo
\begin{itemize}
  \item Clase P
  \item Clase NP
  \item Clase NP-Completo 
  \item Clase NP-Hard
\end{itemize}

Espacio
\begin{itemize}
  \item Clase P-Space
  \item Clase NP-Space
\end{itemize}

\textbf{TIME(f(n)):} Conjunto de problemas de decisión que se pueden resolver mediante una Máquina de Turing Determinista con coste $O(f(n))$.

\textbf{NTIME(f(n)):} Problemas de decisión que se pueden resolver mediante una Máquina de Turing NO Determinista con coste $O(f(n))$.

Lo que una MTD puede resolver en Time($2^n$), una MTND  lo puede resolver en menor tiempo NTIME(n). Pero no al revés, porque no sabemos cuánto más tardara. 

\subsection{Complejidad Temporal}
\subsubsection{Clase P}
Problemas (Lenguajes) decidibles por una MT Determinista (de una sola cinta) en TIEMPO POLINÓMICO. $P= \bigcup_k TIME(n^k)$
\begin{itemize}
  \item Time($n$), Time($n^2$), ..., Time($n^k$)
\end{itemize}

La clase P contiene todos los problemas decibles por un algoritmo que cuyo coste computacional (tiempo) es de orden polinómico con respecto al tamaño de la entrada (n).

''P se corresponde con la clase de problemas que se pueden resolver desde un punto de vista realista/práctico en un computador''

La mayoría de los problemas corrientes pertenecen a la clase P: Ordenación, Búsqueda, Operaciones con vectores.

“La clase P es robusta”
\begin{itemize}
  \item La clase P es invariante (i.e. robusta) para todos los modelos de computación que son polinómicamente equivalentes a la MT Determinista de una sola cinta.
  \item Los problemas en P permanecen en P aunque se cambie el modelo de computación.
\end{itemize}

Ejemplos de problemas en P:
\begin{itemize}
  \item Si existe un camino entre dos puntos en un grafo
  \item Si x e y son coprimos
  \item Si x es primo (AKS lo demuestra)
\end{itemize}

\subsubsection{Clase NP}
Problemas (Lenguajes) decidibles por una MT NO Determinista (de una sola cinta) en TIEMPO POLINÓMICO. $P= \bigcup_k NTIME(n^k)$
\begin{itemize}
  \item Time($n$), Time($n^2$), ..., Time($n^k$)
  \item NTime($n$), NTime($n^2$), ..., NTime($n^k$)
\end{itemize}

Todo P es NP, $P \leq NP$, pero no se ha podido demostrar que sean el mismo conjunto, por ahora $P \subset NP$.

Podemos evitar la fuerza bruta para conseguir soluciones de muchos problemas en tiempo polinómico.

Si un problema no es NP, no es TD.

No se conoce (no que no exista) ... en tiempo polinómico.

Lo primero y más fácil que podemos utilizar para saber si pertenece a NP es hace un verificador, y si este lo hace en tiempo polinómico.

El segundo paso es ver si ese puede resolver con una MT ND.

\subsubsection{Verificabilidad}
Verificador V para un lenguaje L: $L=\{w / V \textit{ acepta } <w,c> \textit{ para alguna palabra }c\}$
$$w \in L \Rightarrow \exists c \in \Sigma^*. V \textit{ acepta } <w, c>$$
\begin{itemize}
  \item \textbf{Verificador:} Es un algo que permite ver si una instancia/cadena cumple el problema o no, pero no lo resuelve. Dice si la solución es buena o no. 
  
  El verificador asociado a un problema debe tener para cada una de las palabras de dicho problema (w) una cadena c llamada certificado que hace que termine el verificador en un estado de aceptación. Para aquella que no sean del lenguaje las rechazara, no habrá un c.
  \item \textbf{Certificador:} Cadena que se asocia a una salida w que recibe el verificador que ayuda a saber si la solución pertenece al lenguaje.
\end{itemize}

NP es la clase de lenguajes que tienen verificadores en tiempo polinómico.

Encontrar un verificador es más sencillo que resolver el problema, este solo dice si es una buena solución o no resuelve el problema.

\subsubsection{P vs. NP}
P es la clase de lenguajes en los que ... (reconocer/aceptar) \dots

NP ... (verificar) ...

\subsubsection{Clase EXPTIME}
$EXPTIME= \bigcup_k TIME(2^{n^k})$

Todos los conjuntos exponenciales. $NP \leq EXPTIME$

\subsubsection{Clase NP-Completo}
Un problema de decisión $p\in NP-Completo$ si y solo si:
\begin{itemize}
  \item $p \in NP$
  \item $\forall q \in NP$ existe una reducción de q en p. Es decir, existe una transformación f de coste polinómico de q en p, tal que $q \leq_p p$
\end{itemize}

Los problemas NPC se pueden reducir entre ellos.

Todos los problemas de NP se pueden reducir en tiempo polinómico en NP-Completo. De esta manera se puede demostrar que un problema pertenece a NP-Completo, si se puede reducir en tiempo polinómico un problema NP-Completo a este.
\begin{itemize}
  \item Tenemos los problemas $p$ y $ q\in NP-Completo$, por lo que si $q \leq_p p$ entonces $p \in NP-Completo$.
  \item Ejem. Como el Buscaminas (Minesweeper) es más difícil que SAT y sabemos que SAT es NP-Completo entonces Buscaminas es NP-Completo.
\end{itemize}

Formas de resolver problemas en tiempo razonable:
\begin{itemize}
  \item Casos particulares: Cuando se conocen casos particulares del problema para los cuales existen soluciones rápidas.
  \item Heurísticas: Un algoritmo que trabaja razonablemente bien en muchos casos. En general son rápidos, pero\dots
  \item Metaheurísticas:
\end{itemize}

3SAT es NP-Completo, pero 2SAT no.

\subsubsection{3SAT $\leq$ K-Clique}
K-Clique sabemos que es NP, vamos a demostrar si es NP-Completo.
$$\textit{3SAT} \leq \textit{K-Clique}$$

Cada literal es un vértice del clique. Con 3 cláusulas tenemos 9 literales, por tanto 9 vértices. $\phi=(x\vee y \vee z) \wedge (x\vee x \vee \overline{y}) \wedge \dots$

Entonces se hacen todas las conexiones posibles siguiendo:
\begin{itemize}
  \item No se conectan literales en la misma cláusula.
  \item Un literal no se conecta con su opuesto. $y e \overline{y}$
\end{itemize}

La solución será aquella que haga True todas las cláusulas, que coincide con un clique en el grafo que hemos generado.

\begin{enumerate}
  \item $w\in \textit{3SAT} \rightarrow f(w)\in \textit{K-Clique}$ y $f(w)$ se puede hacer en tiempo polinómico.
  
  Supongamos que $\phi$ tiene una asignación satisfacible de literales.
  \begin{itemize}
    \item Entonces un literal por lo menos en una cláusula debe tener valor de verdadero.
    \item ...
  \end{itemize}
  \item $f(w)\in \textit{K-Clique} \rightarrow w \in \textit{3SAT}$
  \begin{itemize}
    \item Se hace lo inverso, se construye un 3SAT a partir de un grafo.
    \item No habrá vértices en el grafo que corresponda a la misma cláusula.
  \end{itemize}
\end{enumerate}

\subsubsection{Clase NP-Hard}
Un problema de decisión $h \in \textit{NP-Hard}$ si y solo si:
\begin{itemize}
  \item $\forall q \in NP$, existe una reducción de q en h. Es decir, existe una transformación f de coste polinómico de q en h tal que $q \leq_p h$
\end{itemize}

Los NP-Hard pueden reducir los NP sin ser NP, siendo NTD.

Todos los NPC son NPH, $\textit{NPC}=\textit{NP} \cap \textit{NP-Hard}$, pero no todos los NPH son NPC porque no están en NP.

Se puede saber que un NPH está en NP, si se puede reducir a un NPC.

\tikzset{every picture/.style={line width=0.75pt}}       
\begin{tikzpicture}[x=0.75pt,y=0.75pt,yscale=-1,xscale=1]
\draw  [line width=1.5]  (140,30) -- (290,30) -- (290,100) -- (140,100) -- cycle ;
\draw  [line width=0.75]  (140,70) -- (290,70) -- (290,140) -- (140,140) -- cycle ;
\draw    (210,88.5) -- (227,88.5) ;
\draw [shift={(230,88.5)}, rotate = 180] [fill={rgb, 255:red, 0; green, 0; blue, 0 }  ][line width=0.08]  [draw opacity=0] (8.93,-4.29) -- (0,0) -- (8.93,4.29) -- cycle    ;
\draw    (225,88.5) -- (198,88.5) ;
\draw [shift={(195,88.5)}, rotate = 360] [fill={rgb, 255:red, 0; green, 0; blue, 0 }  ][line width=0.08]  [draw opacity=0] (8.93,-4.29) -- (0,0) -- (8.93,4.29) -- cycle    ;

\draw    (160,110) -- (198.34,52.5) ;
\draw [shift={(200,50)}, rotate = 123.69] [fill={rgb, 255:red, 0; green, 0; blue, 0 }  ][line width=0.08]  [draw opacity=0] (8.93,-4.29) -- (0,0) -- (8.93,4.29) -- cycle    ;
\draw    (270,110) -- (231.66,52.5) ;
\draw [shift={(230,50)}, rotate = 56.31] [fill={rgb, 255:red, 0; green, 0; blue, 0 }  ][line width=0.08]  [draw opacity=0] (8.93,-4.29) -- (0,0) -- (8.93,4.29) -- cycle    ;
\draw    (190,80) -- (207.88,62.12) ;
\draw [shift={(210,60)}, rotate = 135] [fill={rgb, 255:red, 0; green, 0; blue, 0 }  ][line width=0.08]  [draw opacity=0] (8.93,-4.29) -- (0,0) -- (8.93,4.29) -- cycle    ;
\draw    (240,80) -- (222.12,62.12) ;
\draw [shift={(220,60)}, rotate = 45] [fill={rgb, 255:red, 0; green, 0; blue, 0 }  ][line width=0.08]  [draw opacity=0] (8.93,-4.29) -- (0,0) -- (8.93,4.29) -- cycle    ;

\draw (236,79.4) node [anchor=north west][inner sep=0.75pt]    {$q_{3}$};
\draw (177,79.4) node [anchor=north west][inner sep=0.75pt]    {$q_{2}$};
\draw (292,33) node [anchor=north west][inner sep=0.75pt]   [align=left] {NPH};
\draw (292,73) node [anchor=north west][inner sep=0.75pt]   [align=left] {NPC};
\draw (297,111) node [anchor=north west][inner sep=0.75pt]   [align=left] {NP};
\draw (262,112.4) node [anchor=north west][inner sep=0.75pt]    {$q_{4}$};
\draw (152,112.4) node [anchor=north west][inner sep=0.75pt]    {$q_{1}$};
\draw (209,36.15) node [anchor=north west][inner sep=0.75pt]    {$h$};
\end{tikzpicture}

\subsection{Complejidad Espacial}
Se utiliza como otra forma de clasificar problemas, según su dificultad computacional.

El modelo matemático para esta clasificación es la MT.

La complejidad espacial de una MTD M, que para siempre, es una función $f: N \rightarrow N$, tal que $f(n)$ es el máximo número de celdas que M visita para una entrada de longitud n.

Cantidad de espacio (memoria) que se requiere. El espacio a diferencia del tiempo se puede reutilizar.

Se dice que M se ejecutara en un espacio $f(n)$:
\begin{itemize}
  \item Ejem: $<MTD, w>$ MTD se ejecuta en $f(n)=4n$ para $n=2$ $f(2)=8$. Con entrada 2 usará un total de 8 celdas. 
\end{itemize}

\textbf{SPACE(f(n)):} Problemas (Lenguajes) decidibles por una MT Determinista (de una sola cinta) en ESPACIO POLINÓMICO. $PSPACE=\bigcup_k SPACE(n^k)$

\textbf{NSPACE(f(n)):}Problemas (Lenguajes) decidibles por una MT NO Determinista (de una sola cinta) en ESPACIO POLINÓMICO. $NPSPACE=\bigcup_k NSPACE(n^k)$

SAT por ejemplo es $O(n)$ en cuanto a espacio y en tiempo es exponencial.

\subsubsection{Teorema de Savitch (1970)}
$\forall f: N \rightarrow R^+$, donde $f(n)\geq n$. $\textit{NSPACE}(f(n)) \subseteq \textit{SPACE}(f^2(n))$

En otras palabras, cualquier problema decidible en una máquina de Turing no determinista en el espacio $f(n)$ también se puede decidir en una máquina de Turing determinista en el espacio $f^2(n)$.

Corolario: $\textit{PSPACE}(n^4)=\textit{NPSPACE}(n^k)$. Como los PSPACE absorben los NPSPACE resultan ser lo mismo.
\begin{itemize}
  \item $P \subseteq NP$. P incluido en NPC
  \item $P \subseteq \textit{PSPACE}$
  \item $NP \subseteq \textit{NPSPACE}$
  \item $\textit{PSPACE} = \textit{NPSPACE}$
  \item $P \subseteq NP \subseteq \textit{PSPACE}=\textit{NPSPACE} \subseteq \textit{EXPTIME}$
\end{itemize}

\chapter{Modelos de Computación}\label{ch:modelos-de-computación}
Es un modelo que describe la estructura que hace posible computar una salida a partir de una entrada, conocido como modelo, se puede medir la complejidad de un algoritmo que resuelve un problema.

Un modelo recibe como entrada una instancia, aplica un algoritmo y se obtiene una salida, analizando este proceso se determina la complejidad.

Algunos modelos de computación:
\begin{itemize}
    \item Automatas Finitos (AFD)
    \item Automatas de Pila (AP)
    \item Máquinas de Turing (MT)
    \item Random Access Machine (RAM)
    \item Automatas Celulares (AC)
\end{itemize}

\textbf{Turing Completo:} Un modelo es Turing-completo si puede computar (resolver) las mismas funciones (problemas) que una Máquina de Turing.
\begin{itemize}
    \item AF y AP NO son Turing Completos.
    \item RAM, $\lambda$-calculus y Automatas Celulares son Turing Completos.
\end{itemize}

Un problema No-Turing-Decidible aunque se cambie el modelo de computación no lo va a dejar de ser.

La MT es el modelo de modelos. La MT es muy potente en capacidad de resolver problemas, pero muy ineficiente en cómputo, por lo que utilizar otras puede resultar más eficiente, pero no son más capaces.

\section{Automatas Finitos (AF)}
Se definen mediante: ($\Sigma$, Q, f, $q_0$, F)
\begin{itemize}
    \item $\Sigma$: alfabeto de entrada.
    \item Q: conjunto de estados.
    \item f: función de transición.
    \item $q_0$: estado inicial, $q_0 \in Q$.
    \item F: conjunto de estados finales, F $\subset$ Q.
\end{itemize}

Pueden ser Deterministas (AFD) y No Deterministas (AFN). Se puede pasar un AFN a un AFD.

Es No Determinista si:
\begin{itemize}
    \item Para un mismo signo en un estado hay más de una transición.
    \item Transitar con $\lambda$.
    \item ?No hay estado final?
\end{itemize}

El problema de estas máquinas es que no tiene memoria, por lo que no pueden recordar.

\section{Automata de Pila (AP)}
Se definen mediante: ($\Sigma$, $\Gamma$, Q, $A_0$, $q_0$, f, F)
\begin{itemize}
    \item $\Sigma$: alfabeto de entrada.
    \item $\Gamma$: alfabeto de pila (normalmente son mayúsculas).
    \item Q: conjunto de estados.
    \item $A_0$: conjunto de estados de la pila inicial, $A_0 \subset Q$.
    \item $q_0$: estado inicial, $q_0 \in Q$.
    \item f: función de transición.
    \item F: conjunto de estados finales, F $\subset$ Q.
\end{itemize}

Un APN NO se puede pasar a APD.

\section{Generalizaciones de las Máquinas de Turing}
Toda posible generalización de la Máquina de Turing se puede simular con una máquina de Turing: 
\begin{itemize}
    \item Determinista, 
    \item con cinta infinita en ambos sentidos, y \item donde el alfabeto de la cinta es {1,0, b} (b es celda vacía, en blanco)
\end{itemize}

\section{Random Access Machine Model (RAM)}
También conocido como Máquina de Registros.

Modelo abstracto de los computadores digitales. Reflejan más los aspectos de la computación real que las MT.

MT debe recorrer las celdas anteriores a una concreta para leerla, esta, sin embargo, puede acceder directamente a los registros al estar numerados.

Estructura:
\begin{itemize}
    \item $R_0$, $R_1$, ...; Un número ilimitado de registros (celdas). Capaz de almacenar un entero cualquiera.
    \item $I_0$, $I_1$, ...; La entrada es un conjunto finito de datos, no almacenados en los registros.
    \item \textbf{Programa:} Es un conjunto finito de instrucciones, que no están almacenados en los registros. Almacenado en un número finito de registros, y uno adicional, denominado Program Counter (PC).
\end{itemize}

Conjunto de posibles instrucciones (cada operación aumenta el PC, excepto el salto condicional que llega a un determinado valor):
\begin{itemize}
    \item $R_i \leftarrow R_j$ Asignación
    \item $<R_i> \leftarrow R_j$ y $R_i \leftarrow <R_j>$ Direccionamiento indirecto ($<R_i>$ valor almacenado en la celda con dirección almacenada en $R_i$)
    \item $R_i \leftarrow R_j + R_k$ Suma
    \item $R_i \leftarrow R_j - R_k$ Resta
    \item $R_i \leftarrow R_j \textit{ bool } R_k$ Operación booleana
    \item If $R_i$ comp $R_j$ label 1 else label 2 Salto condicional, para comp=comparación
\end{itemize}

Hay diferentes tipos de arquitecturas (como pasa con los tipos de MT) con las que se puede implementar: Counter Machine, Pointer Machine, Random Access Machine, Random Access Stored Program Machine.

\subsection{Random Access Stored Program (RASP)}
Es una máquina abstracta empleada para implementar algoritmos.

Se diferencia de la RAM, en que el programa está almacenado en los propios registros de la máquina. Los datos y las instrucciones están en la misma memoria (memoria de acceso aleatorio)
\begin{itemize}
    \item RAM $\rightarrow$ Arquitectura Von Harvard $\rightarrow$ Máquina de Turing.
    \item RASP $\rightarrow$ Arquitectura Von Neumann $\rightarrow$ Máquina de Turing Universal.
\end{itemize}

\section{Automatas Celulares}
Sucesión de sitios que contienen un valor numérico, que está determinado a seguir unas reglas que hacen evolucionar estos valores a lo largo del tiempo.

Es un modelo de un ''mundo'' con una física muy sencilla. El mundo está dividido en celdas (entramado de celdas).

El Automata Celular está gobernado por reglas que rigen la evolución del sistema en el tiempo.
\begin{itemize}
    \item El tiempo se considera Discreto (steps).
    \item Las reglas para el siguiente paso en el tiempo (t+1) se basan en el estado del autómata en el instante actual (t).
    \item El comportamiento de una celda se puede representar (computar) mediante un Automata Finito.
\end{itemize}

\textbf{AC Deterministas:} Las reglas no tienen componentes de aleatoriedad, dado el mismo estado inicial, siempre se obtendrán el mismo estado final.

\textbf{AC No Deterministas:} Sigue una aleatoriedad controlada, pero aleatoria.

Aplicaciones: Generadores de contenido multimedia, Generadores de números aleatorios, Desarrollo de Computadores paralelos y Supercomputación.

\subsection{Variedades}
Pueden diferir en:
\begin{itemize}
    \item Dimensión del entorno d=1,2,3...
    \item Formación de celdas: triangular, cuadrado, rectangular, pentagonal, hexagonal, etc.
    \item Forma y tamaño del entorno: triangular, cuadrado, rectangular; finita, infinita.
    \item Tipo se automata finito asociado a cada celda: determinista, probabilístico.
    \item Definición de los vecinos de cada automata
    \begin{itemize}
        \item Puede indicarse con un número r (distancia), en función del cual obtenemos el número de vecinos que rodean a una celda concreta.
        \item En AC bidimensionales podemos considerar: 
        \begin{itemize}
            \item Vecindad de Von Neumann (arriba, abajo, izquierda y derecha)
            \item Vecindad de Moore (todas las direcciones, incluyendo las diagonales).
        \end{itemize} 
        \item Número de celdas que rodean a una celda concreta central: $(2r+1)^2-1$
    \end{itemize}
\end{itemize}

\subsection{Clasificación}
\begin{itemize}
    \item Clase I: Desde casi cualquier estado inicial convergen siempre a una misma configuración final.
    \item Clase II: Genera un patrón simple, donde hay una estructura que se repite.
    \item Clase III: Generan patrones aleatorios.
    \item Clase IV: Generan estructuras locales que son capaces de sobrevivir durante mucho tiempo (mantenerse iguales).
\end{itemize}

Un ejemplo de este tipo de autómata, es el Juego de la Vida de John H. Conway, en el que la configuración va evolucionando de manera que parece una población.
\begin{itemize}
    \item Cada célula está vida o está muerta (estados binarios).
    \item Tiene una vecindad de r=1, de tipo Moore (todas las adyacentes).
    \item Se van formando diferentes patrones.
    \item Regla: 
    \begin{itemize}
        \item Muere si hay menos de 2 o más de 3 células vivas alrededor.
        \item Vivos si hay exactamente 2 o 3 células vivas alrededor.
        \item Nace si es un 0 y hay exactamente 3 células vivas alrededor.
    \end{itemize}   
\end{itemize}

\subsection{AC Unidimensional}
Dimensión del AC: Dimensiones del entramado que forman el conjunto de celdas.

En una dimensión, hay 3 configuraciones:
\begin{itemize}
    \item \textbf{Secuencia finita:} Hay posibles estados finitos.
    \item \textbf{Anillo:} Siempre vuelve al mismo estado.
    \item \textbf{Secuencia infinita:} El cambio siempre es diferente.
\end{itemize}

Las Reglas dependen del estado que tenga la celda $x_i$ en el instante anterior y del estado del vecindario (izq. y dch.) en el instante anterior.

Rule 50 y 30, ejemplo de reglas que con vecindario de las celdas adyacentes (3 valores) define que valor tomara la siguiente, la distribución será el valor en binario 50 o 30 escritos de menos a más significativo de arriba a abajo.

\section{L-Systems}
Propuesto por Aristid Lindemayer en 1968.

Es un tipo de lenguaje formal generado por una gramática. Es Turing-Completo.

Aplicación: Simulación de procesos biológicos (plantas, hojas, etc.), Turtle interpretation, Patrones bidimensionales, Formalización de curvas fractales.

\textbf{Definición formal:} ($\Sigma$, w, P)
\begin{itemize}
    \item $\Sigma$: Alfabeto de símbolos.
    \item w: Axioma, $w \in \Sigma^*$.
    \item P: Conjunto finito de reglas del tipo u::=v donde $u \in \Sigma^*$ y $v \in \Sigma^*$.
\end{itemize}

\textbf{Lenguaje generado por un L-system:} $\{x | u \in \Sigma^+, w \rightarrow^* x\}$. Las reglas se aplican en Paralelo.

\textbf{Tipos de Gramáticas:}
\begin{itemize}
    \item \textbf{Sensibles al contexto (1L):} La parte derecha de la producción sea igual o más larga que la izquierda, y esta puede tener más de un símbolo. $xay::=xvy$
    \item \textbf{Independientes del contexto (0L):} $V \rightarrow w$. Donde V es un símbolo no terminal y w es una cadena de terminales y/o no terminales
    \item \textbf{Deterministas (DL):} solo existe una regla posible que ser aplicada.
    \item \textbf{Propagativos (PL):} la parte derecha no puede ser más corta que la izquierda.
    \item \textbf{Con Tablas (TL):} en cada paso de la derivación se aplican distintos conjuntos de reglas.
    \item \textbf{Con Extensiones (EL):} solo pertenecen al lenguaje las palabras formadas por un subconjunto del alfabeto.
\end{itemize}

\subsection{L-Systems vs. Gramáticas de Chomsky}
\textbf{Gramáticas de Chomsky:} Las reglas se aplican de una en una, Derivación Secuencial.

\textbf{L-Systems:} Todas las derivaciones se realizan Simultáneamente, Derivación en paralelo. Se aplican las reglas para todos los símbolos (distintos o no) de la palabra.

\section{Máquinas de Turing}
Se definen mediante: (Q, $\Gamma$, b, $\Sigma$, $\delta$, $q_0$, F)
\begin{itemize}
    \item Q: conjunto de estados.
    \item $\Gamma$: alfabeto de cinta.
    \item b: símbolo en blanco, b $\in$ $\Gamma$.
    \item $\Sigma$: alfabeto de entrada.
    \item $\delta$: función de transición.
    \item $q_0$: estado inicial, $q_0 \in Q$.
    \item F: conjunto de estados finales, F $\subset$ Q.
\end{itemize}

\subsection{Origen}
Formulada por Turing como un medio para formalizar el proceso sistemático de resolver problemas de un matemático.

En sentido más amplio: Para representar cualquier Algoritmo, por precaución hablamos de Procedimientos Efectivos para resolver un problema.

Actualmente, es una máquina evolucionada en base a contribuciones de Post, Turing y otros.

\subsection{Utilidad}
Turing se sirvió de ella para desarrollar diversas cuestiones: 
\begin{itemize}
    \item Máquina de Turing Universal.
    \item Problema de la parada (demostración computacional) en equivalencia con el Entscheidungsproblem de Gödel.
    \item Modelo Universal de Computación.
\end{itemize}

\subsection{Modelo de Computación Universal}
Las Máquinas de Turing definen un Modelo de Computación que es la base de los Lenguajes Imperativos y Procedimentales.

Suposiciones previas:
\begin{itemize}
    \item MT de cinta infinita
    \item Alfabeto$=\{b,1\}$ siendo b o $\#$ representan celda en blanco
    \item Alfabeto auxiliar con símbolos de marcado $\{\$, X, etc\}$
    \item Se representan Números Naturales en base 1
\end{itemize}

\subsubsection{Aritmética de Peano} 
Conjunto axiomático que sirve para definir los Números naturales
\begin{itemize}
    \item R0 - Objeto inicial. 0 es un objeto especial y es un Número Natural.
    \item R1 - Sucesor. Para cada n NN existe otro NN $succ(n)$.
    \item R2 - Predecesor. Si a es $succ(b)$ entonces b es $pred(a)$.
    
    Ojo con $pred(k=1)$, no está definido si no incluimos 0
    \item R3 - Única. No hay dos NN diferentes con el mismo sucesor.
    \item R4 - Igualdad. Comparar a y b NN en cuanto a igualdad.
    \begin{itemize}
        \item Reflexiva $Eq(a,a)$
        \item Simétrica $Eq(a,b) \Leftrightarrow Eq(b,a)$
        \item Transitiva $Eq(a,b), Eq(b,c) \Leftrightarrow Eq(a,c)$
    \end{itemize}
    \item R5 - Inducción. Predicado $P(n)$ es cierto $\forall n$ si cumple: 
    
    $P(0)$ cierta.
    
    $P(n)$ cierta para algún n y se puede demostrar que $P(succ(n))$ es cierta.
\end{itemize}

\subsection{Eficiencia}
No, no busca la eficiencia.

La MT es un modelo de computación abstracta.

Igual sabe hacer cosas que podríamos calcular de otra manera. No obstante, vamos a estudiar la eficiencia $\Rightarrow$ Coste Computacional.
\pagebreak

\subsection{Coste Computacional con MT}
Uso para el Análisis del Coste Computacional de Algoritmos.

\subsubsection{Ventajas}
\begin{itemize}
    \item Propone un modelo único, al contrario que con el uso de un Lenguaje X y una CPU Y
    \item No hay dudas sobre el coste del Paso Base
\end{itemize}

\subsubsection{Desventaja}
\begin{itemize}
    \item Es muy pesado para algoritmos complejos $\Rightarrow$ Limitaremos el tamaño de las MT
\end{itemize}

\subsubsection{Propósito}
\begin{itemize}
    \item Relacionar conceptos vistos en la asignatura.
    \item En base a un modelo muy sencillo
    \item Abstraer de construcciones de muy alto nivel
\end{itemize}

\subsubsection{Marco necesario para el Análisis}
\begin{itemize}
    \item MT de cinta infinita hacia ambos lados
    \item Alfabeto determinado, en principio con pocos símbolos, pero pueden ser más
    \item Una instancia de un Problema se describirá con una cinta de Entrada
    \item El número de símbolos de la Entrada define el Tamaño de la Instancia
    \item Coste de Paso Base: Una Transición con su Lectura/Escritura y Desplazamiento
\end{itemize}

\subsubsection{Coste computacional empírico: Aplicando diferencias finitas}
\begin{table}[H]
    \begin{tabular}{|l|l|l|l|l|l|l|}
    \hline
    Tamaño n                  & \textbf{1} & \textbf{3} & \textbf{5} & \textbf{7} & \textbf{9} & \textbf{11} \\ \hline
    \textbf{Pasos}            & 2          & 9          & 20         & 35         & 54         & 77          \\ \hline
    \textbf{A: T(n) - T(n-2)} &            & 7          & 11         & 15         & 19         & 23          \\ \hline
    \textbf{B: A(n) - A(n-2)} &            &            & 4          & 4          & 4          & 4           \\ \hline
    \textbf{C: B(n) - B(n-2)} &            &            &            & 0          & 0          & 0           \\ \hline
    \end{tabular}
    \caption{Diferencias finitas}
\end{table}

Si $T(n)$ \textbf{exponencial}, las diferencias finitas no se anulan, pero aparecen comportamientos particulares.

Si $T(n)$ \textbf{polinómico} de orden k, las diferencias finitas de orden $k=cte$. En este caso orden 2.

Sabiendo el orden se plantea un polinomio genérico de ese orden: $T(n)=xn^2+bn+c$.

Se resuelve la ecuación y ya tenemos la complejidad.

$$a+b+c=2;\; 9a+3b+c=9;\; 25a+5b+c=20 \Rightarrow a=\frac{1}{2},\; b=\frac{3}{2},\; c=0$$

$$T(n)=\frac{1}{2}n^2+\frac{3}{2}n$$

\subsubsection{Coste computacional empírico: Recursivamente}
Se plantean las observaciones y analizan los patrones de los datos.

\begin{itemize}
    \item $n=1 \Rightarrow 2$ pasos
    \item $n=3 \Rightarrow 4$ pasos a derecha $+ 3$ a izquierda $+ T(1)$
    \item $n=5 \Rightarrow 6$ pasos a derecha $+ 5$ a izquierda $+ T(3)$
    \item $n \Rightarrow n+1$ pasos a derecha $+ n$ a izquierda $+ T(n-2) = \Sigma i \Rightarrow T(n)=\frac{(n+1)(n+2)}{2}-1=\frac{1}{2}n^2+\frac{3}{2}n$ 
\end{itemize}

\subsection{Variantes del Modelo de Máquina de Turing}
Abordamos:
\begin{itemize}
    \item Máquinas de Turing de dos o más cintas (multicinta)
    \item Máquinas de Turing No Deterministas
    \item Automata de Dos Pilas
    \item Máquinas de Turing 2D
    \item Aplicaciones
    \item Costes Computacionales
\end{itemize}

\subsubsection{Máquinas de Turing de Dos Cintas}
La misma estructura que una MT convencional, pero:
\begin{itemize}
    \item Tiene dos cintas y sus respectivos cabezales
    \item Pueden operar de forma independiente
\end{itemize}

Implica cambios en el funcionamiento, las transiciones se definen:
\begin{itemize}
    \item Cambio de estado
    \item Para cada cabezal:
    \begin{itemize}
        \item Una Lectura
        \item Una Escritura
        \item Un Desplazamiento
    \end{itemize}
\end{itemize}

Las operaciones de cada cabezal son independientes entre sí, pero están vinculados en cada transición.

\subsubsection{Máquinas de Turing No Deterministas}
La misma estructura que una MT convencional, pero
\begin{itemize}
    \item Admiten transiciones No Deterministas, es decir, para un mismo símbolo leído define al menos dos transiciones que:
    \begin{itemize}
        \item Escriben símbolos diferentes
        \item Desplazan el cabezal de distinta forma
        \item Transitan a estados diferentes
    \end{itemize}
\end{itemize}

Cuando se encuentra una transición ND, entendemos que la MT bifurca en dos o más instancias de la misma que adoptaran estados diferentes.

Si alguna de las instancias alcanza un Estado Final $\Rightarrow$ MT acepta

\end{document}